% !TeX root = ../papers/nc-p3.tex
\section{全標準中間辺と可約割線六次への縮約}
\label{sec:integral-middle-secants}
\begin{theorem}\label{thm:integral-middle-secants}
二十次の整な一般被約余因子曲線の二次曲面射影では次が必要である。
\begin{enumerate}
\item 滑らかな二次曲面のどちらの標準中間辺も有限被約であり、
同じ標準中間座標への折り返しはどちらの側でも三次交叉が四未満となる。
\item 錐の場合、標準中間の三次・五次表示の割線平面六次は整であってはならない。
滑らかな場合、二つの標準中間の四次・四次表示の割線平面六次は\emph{両方とも}整であってはならない、
すなわちどちらの側にも整な六次は許されない。
これらは全中間座標の選択に対する正則性の必要条件である。
\end{enumerate}
\end{theorem}
\begin{proof}[滑らかな二次曲面での有限性と折り返し]
各母線で $V$ の二次制限が全射になることは錐の場合と同じ入力直線の排除で従う。
定理\ref{thm:quadric-ruling-edge-obstruction}の二つの二項四次同定による $T$ を使う。
$y$ が可約な平面切断なら中間式は $l\,T(u m)=0$（$l,u$ は非零一次式）となり解はない。
既存の三点繊維の平面も可約ではない。各母線に沿って埋込みだから二つの底への射はともに
長さ三繊維上の閉埋込みであり、各母線との交長は一以下である。
より直接には、いずれの有限長繊維でも二つの一次因子の積が零なら、
長さはそれぞれの因子による商長の和以下である。
滑らかな平面切断上では $S_{23}/yS_{12}$（または $S_{32}/yS_{21}$）が二項五次式、
$V$ の制限 $W$ は二項四次系である。従って中間解は $(S_5/WS_1)^*$。
定理\ref{thm:cone-fold-plucker}の四次系の証明がそのまま使え、解は実際の長さ三繊維の平面だけで一次元である。
他方、固定 $m$ に対する平面が二次元以上の線形空間の解を持てば無限個の $y$ が出る。
従って両射影は有限であり、両標準辺は二十の既存の相異なる印で尽きる。
$m$ が相異なるのは有限線形繊維より従う。

二つの四次評価曲線は交わらないことも一般性を要しない。
母線 $z$ の一点で消える一次式を $l_t$、消えない一次式を $u_t$ とし、
そこで恒等零な $v_z\in V$ を取る。他方の母線で同様に $l_s,u_s,v_w$ を取る。
四次汎関数の評価点を $S_{22}/V$ で表すと $u_t v_z/l_t, u_s v_w/l_s$ に比例する。
実際 $v_z/l_t\in S_{21}$ は $l_t$ 倍が $V$ に入り、評価列の双対で一点評価の直線を生成する。
もし二点が一致すれば $l_tl_s$ を掛けて $V\otimes S_{11}\to S_{33}$ に関係を得る。
$u_tl_s,u_sl_t,l_tl_s$ は独立なので非零関係であり、積同型に反する。

$m,y,m$ の折り返しの正の三次交叉を計算する。$u_0,u_1$ を $t$ 側、$s_0,s_1$ を $s$ 側の一次基底とする。
対称な交叉元は線形写像 $L:S_t^3\to(S_s^5)^*$ で
$L(p)(s_jT(u_i q))=L(q)(s_jT(u_i p))$ を満たすものとなる。
各 $j$ で二項係数を比較すると $L(p)(s_j l)=\gamma_j(p T^{-1}l)$、
$\gamma_j\in(S_t^7)^*$ が一意に定まる（$L_{a,b+i}=L_{b,a+i}$ は添字和ごとの一致である）。
二つの表示を $s_0s_1 r$ で一致させる条件から、交叉は次の写像の余核の双対になる。
\[
   S_t^3\otimes S_s^3\longrightarrow (S_t^7)^2,\qquad
   p\otimes r\longmapsto
   (p T^{-1}(s_1r),-p T^{-1}(s_0r)).
\]
底 $\mathbb P^1_t$ で $r$ から二つの四次式への写像は各点で全射である。
階数落ちは、変換された四次汎関数の二つの一次収縮が従属、すなわちそれが純四次評価であることを意味し、
上の二曲線の交点となる。従って核束は $\mathcal O(-a)\oplus\mathcal O(-(8-a))$、
$a\le4$ と置ける。定数核は $T^{-1}(s_1r)=0$ から零なので $a\ge1$。
三捻り後の $H^1$ は $4-a\le3$ 次元である。有限辺での反転三次交叉は対称部に入る
（命題\ref{prop:finite-flip-obstructions}）ので正則な折り返しはない。逆の $T$ でも同じである。
\end{proof}

\begin{proof}[割線の階数と正規化]
錐では $(f_m,g_m)=(m,Dm)$、滑らかな場合には $(T(u_0m),T(u_1m))$ とする。
どちらも一次に $m$ に依存し、各底点でその二行の評価は独立である。
後者では二つの四次式は独立、前者では $m$ 自体が非零であり、共通根は有限である。
三次共通因子を持つ入力は標準辺の二十印である。
四次対なら二式の一次シジジーが一意で、このシジジーを $y$ の四係数とする。
三次・五次なら $g_m/m$ の二次式が $y$ を与える。

二次因子 $q$ への二式の制限を合わせた四行束写像は全ての $q$ で階数三以上である。
さもなくば核中の入力直線上で二式を $q$ で割った二式の結果式を取る。
残る次数は $(1,3)$ または $(2,2)$ で、入力直線上の四次式となる。
恒等零なら無限個の印が出る。そうでなければ長さ四の零スキームが印のスキームに入る。
実際前者は残余一次式による割り算、後者は共通因子ちょうど三のとき一意な一次シジジーを解くことで、
結果式の零上で標準辺の $y$ をスキームとして復元する。
後者で局所的にこれを見るには残余二次対の二次Bezout行列の階数が一であることを使えばよい。
単元成分で核の一次因子を解き、行列式で割った環上で両二次式がこの因子の倍となる。
共通因子四は独立性で起きず、前者の残余一次式も恒等零にならない。
標準辺は被約なのでこの長さ四も相異なる四印である。実際の中間配置 $U$ と標準 $m$ の二乗評価空間は
$y$ の有限辺の完全跡の直交性により標識点上の単元倍で一致する。従って標準の四点が直線上なら、
対応する $U$ の四点でも二次評価階数は三以下である。相異なる非共線四点の二次評価は全射である
（各一点の外の三点を二平面で覆ってその一点を避けられる）から $U$ にも四点直線ができる。
実際の四辺窓に対する系\ref{cor:finite-no-four-on-line}に反する。

割線六次 $P$ をこの四行写像の行列式とする。階数三なら $q$ 上の $m$ は一意に解け、
$P$ は入射スキームそのものである。$P$ が整の場合を考える。
その入力束 $H$ の次数は二十一である。実際
$E_j=(S_j\otimes\mathcal O)/(S_{j-2}\otimes\mathcal O(-1))$ とし、
$\mathbb P^2\times\mathbb P^3$ で $E_{\deg f}\oplus E_{\deg g}$ の入力一捻りの零点類を取る。
入射は余次元四であり、$c_2(E_{\deg f}\oplus E_{\deg g})=21h^2$ が入力次数を与える。
共通因子三の外では $q$ は唯一なので入力像 $C_m$ は次数二十一の整曲線。
正規化母線束 $S$ の結果式曲面への有限射は、二共通根（二重根も含む）が曲線に限られることから
双有理である。像は八次、一般平面切断の正規化は各母線の平面との交点により有理である。
一般平面で $C_m$ との二十一交点は全て特異である。二つの別の逆像があるか、
二重共通根による非被約繊維を持つためである。平面八次の種数差二十一を使い切るため、
各点の正規化余長は一、他の特異曲線はない。従って一般導手は被約である
（一般横断平面への有限双対の制限、または一般曲線点での同じ余長の計算）。
導手が $C_m$ のイデアルそのものになることは前定理の深さの証明による。

$P$ 上の残余二式の結果式は $Z\in|4(H-h)|$ で、長さは六十である。
二十印への逆像も、固定三次式の長さ二部分因子により各三の商長を持つ。
$\mathcal O_P(-Z)\subset I_M\mathcal O_P$ は、四次対では前定理のBezout小行列の恒等式による。
三次・五次でも $M$ 上では $g_m=m a_m$ とスキームとして書ける
（$a_m$ は標準辺で唯一に決まる二次式）。$q$ を割った残余二式は従って前者が後者を割り、
その結果式は $I_M\mathcal O_P$ に入る。商長から両場合で等号である。

前定理の特異平面六次上の次数六の可逆層の議論を適用する。
非恒等根の行列式の $P$ への制限は $8H-3Z$ で次数が負、余因子は $6H-2Z$ に入る。
従ってこの根の階数は $C_m$ 上で一般に二以下となる。導手の交点と有限双対の同じ証明で
行列式は結果式八次に一致し、余核は $\nu_*\pi^*\mathcal O(3)$ である。
錐の場合これは $m(z)=0$ という一次関係を持ち、非恒等根は存在しない。
従って中間はともに標準 $m$ に限られ、前節の折り返し排除に反する。
滑らかな場合はこの根が既存の四次Bezout行列であり、合同を除き一意である。
すなわち二つの中間は標準 $m,n$ だけから選ぶ。異なる両方を取れば
定理\ref{thm:quadric-ruling-edge-obstruction}、同じものなら上の三次交叉に反する。
任意の実際の中間の平方がこの標準平方に揃うことには、有限な $y$ 隣接辺の完全跡の直交性を使えばよい。
\end{proof}
\path{computations/verify_integral_middle_secants.py} は二つの束の次数、全分裂型の折り返し上界と
有理標本での二つの二項乗法行列を検算する。
